% Part: turing-machines
% Chapter: undecidability
% Section: universal-tm

\documentclass[../../../include/open-logic-section]{subfiles}

\begin{document}

\olfileid[tr]{tur}{und}{uni}
\olsection{Evrensel Turing Makineleri}

\olref[enu]{sec} içinde her Turing makinesinin sonlu bir tam sayı dizisiyle
nasıl betimlenebildiğini tartıştık. Bu dizi Turing makinesinin durumlarını,
alfabesini, başlangıç durumunu ve yönergelerini kodlar. Ayrıca bütün bu
betimler kümesinin !!{enumerable} olduğuna dikkat çektik. Böyle betimler
kümesi !!{denumerable} olduğuna göre $\Nat$'ten bu betimlere giden
!!a{surjective} fonksiyon vardır. Böyle !!a{surjective} fonksiyon, örneğin
Cantor'un zikzak yöntemiyle elde edilebilir. Bu bize Turing makinelerinin
(betimlerinin) tümünü numaralandırmanın bir yolunu verir. Böyle bir
numaralandırmayı sabitlersek artık $1$., $2$., \dots, $e$. Turing
makinesinden söz etmek anlamlı olur. Bu sayılara \emph{indis} denir.

\begin{defn}
$M$ (sabit numaralandırmamızdaki) $e$. Turing makinesiyse $e$'nin $M$'nin bir
\emph{indisi} olduğunu söyleriz. $e$. Turing makinesini $M_e$ ile gösteririz.
\end{defn}

Bir makinenin birden fazla indisi olabilir; örneğin $M$'nin iki betimi,
yönergelerini listelediğimiz sıra bakımından farklı olabilir ve bu farklı
betimlerin indisleri de farklı olur.

Önemli olarak, Turing makinesi betimlerinin numaralandırmasını, indisinden
$M$'nin betimini ve $M$ makinesinin betiminden onun bir indisini etkili
biçimde hesaplayabileceğimiz bir biçimde vermek mümkündür. Church--Turing
tezine göre, indisi $e$ olan Turing makinesinin betimini geri elde edip karşılık
gelen betimi çıktı olarak şeridine yazan bir Turing makinesi bulmak da
mümkündür. Betim, $\TMstroke$ bloklarından oluşan bir dizi olurdu (bu bloklar
$M_e$'yi betimleyen dizideki pozitif tam sayıları gösterir).

Bunu bildiğimize göre şu soruları sormak doğaldır: Turing makinesi
indislerinin hangi fonksiyonları Turing makinelerince hesaplanabilir? Turing
makinesi indislerinin hangi özelliklerine Turing makinelerince karar
verilebilir? Bir örnek verelim: $e$ indisini, indisi $e$ olan Turing
makinesinin durum sayısına götüren fonksiyon bir Turing makinesince
hesaplanabilir. Böyle bir Turing makinesi şunu yapardı: Şeritte $e$ tane
$\TMstroke$ simgesinden oluşan tek bir blokla başlatıldığında önce $e$'yi
betimine çözerdi. Betim artık şeritte bir $\TMstroke$ blokları dizisiyle
gösterilir. Bu dizinin ilk !!{element}, durumların sayısıdır. Dolayısıyla
artık yapılması gereken tek şey ilk $\TMstroke$ bloğu dışındaki her şeyi
silmek ve ardından durmaktır.

Dikkate değer bir sonuç şudur:

\begin{thm}\ollabel{thm:universal-tm} $\tuple{e,n}$ girdisinde
  başlatıldığında
  \begin{enumerate}
    \item ancak ve ancak $M_e$, $n$ girdisinde durursa duran ve
    \item $M_e$, $m$ çıktısıyla duruyorsa $U$ da aynı çıktıyla duran
  \end{enumerate}
  bir \emph{evrensel Turing makinesi}~$U$ vardır. Böylece $U$, $M_e$ makinesi
  $n$ girdisinde $m$ çıktısıyla duruyorsa $f(e,n)=m$, aksi durumda tanımsız
  olan $f\colon\Nat\times\Nat\pto\Nat$ fonksiyonunu hesaplar.
\end{thm}

\begin{proof}
  Son derece karmaşık bir makine olduğundan $U$'yu gerçekten üretmek
  neredeyse olanaksızdır. Fakat nasıl çalıştığını ana çizgileriyle betimleyip
  Church--Turing tezine başvurabiliriz. Başlangıçta $U$'nun şeridinde $e$
  tane $\TMstroke$ simgesinden oluşan bir bloğu, $n$ tane $\TMstroke$
  simgesinden oluşan bir blok izler. Önce $e$
  indisini $n$ girdisinin sağında ``çözer''. Böylece $M_e$ makinesinin
  yönergelerini betimleyen bir sayılar listesi (yani araları $\TMblank$
  simgeleriyle ayrılmış $\TMstroke$ blokları) elde edilir. Ardından $U$,
  $M_e$'nin başlangıç durumunun sayısını ve $1$ sayısını $M_e$ betiminin
  sağındaki şeride yazar. (Bunlar da tekli gösterimde $\TMstroke$ blokları
  olarak gösterilir.) Sonra girdiyi ($n$ tane $\TMstroke$ bloğunu) sağa
  kopyalar; ancak her $\TMstroke$ simgesini üç $\TMstroke$ simgesinden oluşan
  bir blokla değiştirir (anımsayın, $\TMstroke$ simgesinin numarası $3$,
  $\TMendtape$ simgesinin numarası $1$ ve $\TMblank$ simgesinin numarası
  $2$'dir). $U$'nun şeridinde $\TMblank$ simgeleriyle ayrılmış bu bloklar
  dizisinin sol ucuna, $\TMendtape$ kodu olan tek bir $\TMstroke$ yazar.

  Artık $U$'nun şeridinde şunlar bulunur: $e$ indisi, $n$ sayısı, başlangıç
  durumunun kod numarası (``mevcut durum''), ilk kafa konumu~$1$'in sayısı
  (``mevcut kafa konumu'') ve ``şeridin'' ilk içeriği ($M_e$ simgelerinin kod
  numaralarını gösteren ve araları $\TMblank$ simgeleriyle ayrılmış bir
  $\TMstroke$ blokları dizisi, yani ``simgeler'').

  Şimdi şu işlemleri yaparak $M_e$'nin $n$ girdisinde başlatıldığında
  yapacaklarını benzetir:
  \begin{enumerate}
    \item ``Mevcut kafa konumunun'' sayısı~$k$'yi bulun (başlangıçta bu
    sayı~$1$'dir),
    \item Oradaki ``simgenin'' ne olduğunu görmek için ``şeritteki'' $k$.
    bloğa gidin,
    \item\ollabel{find-inst}%
    mevcut ``durum'' ve ``simgeyle'' eşleşen yönergeyi bulun,
    \item ``Şeritteki'' $k$. bloğa geri dönüp oradaki ``simgeyi'', $M_e$'nin
    yazacağı simgenin kod numarasıyla değiştirin,
    \item Kafayı mevcut ``durumu'' kaydettiği yere götürüp oradaki sayıyı
    yeni durumun numarasıyla değiştirin,
    \item ``Şerit konumunu'' kaydettiği yere gidip yönerge sırasıyla sola ya
    da sağa gitmeyi söylüyorsa bir~$\TMstroke$ silin ya da
    bir~$\TMstroke$ ekleyin.
    \item Yineleyin.\footnote{Burada bazı ince güçlükleri atlıyoruz. Örneğin
    $U$, ``mevcut kafa konumunu'' izlediği sayacı artırırken ek yere gereksinim
    duyabilir; bu durumda bütün ``şeridi'' sağa kaydırması gerekir.}
  \end{enumerate}
  $M_e$, $n$ girdisinde başlatıldığında hiç durmazsa $U$ da hiç durmaz;
  dolayısıyla çıktısı tanımsızdır.

  \olref{find-inst} adımında $M_e$ betiminin mevcut ``durum''/``simge'' çifti
  için bir yönerge içermediği ortaya çıkarsa $M_e$ dururdu. Bu olduğunda $U$,
  şeridinin ``şeridin'' solunda kalan bölümünü siler. $M_e$'nin şeridindeki
  bir~$\TMstroke$ simgesini gösteren üçlü $\TMstroke$ bloklarının her biri
  için kendi şeridinin sol ucuna bir $\TMstroke$ yazar ve ``şeridi'' art arda
  siler. İşlem bittiğinde $U$'nun şeridinde uzunluğu $m$ olan tek bir
  $\TMstroke$ bloğu bulunur.
  
  $U$, ``şeritte'' üç $\TMstroke$ simgesinden oluşan bir blok dışında bir
  şeyle karşılaşırsa hemen durur. Bu durumda $U$'nun şeridi tek bir
  $\TMstroke$ bloğu içermediğinden çıktısı doğal bir sayı değildir; yani bu
  durumda $f(e,n)$ tanımsızdır.
\end{proof}

\end{document}
